\expandafter\ifx\csname ifdraft\endcsname\relax
 \documentclass{jsarticle}
 \begin{document}
\fi
%%%%%%%%%%%%%%%%%%%%%%%%%%%%%%%%%%%%%%%%%%%

\def\Path{./5_Metaheuristic}

%%%%%%%%%%%%%%%%%%%%%%%%%%%%
\chapter{脳波データに関する理論と手法}

%%%%%%%%%%%%%%%%%%%%%%%%%%%
本章では， 本研究で用いる脳波データから脳状態を取得するために必要な理論と手法について説明する．
具体的には， 注意集中度や課題関与度を反映する指標として用いられる．まず，これまでの研究で提案されてきたEngagement Index（EI），Frontal Alpha/Beta Asymmetry（FAA），Sample Entropyの3つの指標について解説する．
次に，片山研究室H30年度卒業生の平野の研究によって考案された脳波指標について説明する．
%%%%%%%%%%%%%%%%%%%%%%%%%%%%

\section{Engagement Index（EI）に基づく脳波指標}

\subsection{Engagement Indexの概念}

Engagement Index（以下，EI）は，脳波（EEG）信号に含まれる周波数帯域成分の比率を用いて，
被験者の\textbf{注意集中度}や\textbf{課題関与度}を定量化するために提案された指標である．
EIは，主に$\alpha$波，$\beta$波，$\theta$波の相対的なパワー分布に基づいて算出される．
従来研究において，$\beta$波は認知的負荷や能動的思考と関連し，
$\alpha$波および$\theta$波はリラックス状態や注意低下と関係していることが示されている．
この指標を用いた代表的な研究として，駐車場巡回作業中の人間の集中度を脳波計測によって定量的に解明した研究が挙げられる．
この知見に基づき，EIは「覚醒・集中を促進する成分」と
「非集中状態を示す成分」の比として定義される．

\subsection*{認知的負荷と脳波指標}

認知的負荷（cognitive load）とは，
課題遂行時にワーキングメモリや注意資源が
どの程度消費されているかを表す概念である．
認知的負荷が高い状態では，
外部刺激の処理や内部情報の保持・操作が増加し，
前頭部を中心とした脳活動が活性化することが知られている．

脳波においては，
β帯域の増加および
α帯域の抑制が
認知的負荷の上昇と関連づけられており，
これらの周波数特性を組み合わせた指標が
課題集中度の評価に用いられてきた．

\subsection{Engagement Index（EI）ベース指標}

本研究では，
脳波の周波数帯域成分を用いた心理状態評価手法として，
Popeらにより提案されたEngagement Index（EI）を基礎とする指標を使用する．
EIは，
課題遂行時の注意集中や認知的負荷を反映する指標として，
ヒューマンインタフェースやBCI研究において広く用いられている．

\subsection{Engagement Indexの定義}

Engagement Indexは，
β帯域活動と，
α帯域およびθ帯域活動の比として，
次式で定義される．

\begin{equation}
EI = \frac{\beta}{\alpha + \theta}
\end{equation}

β帯域は，
覚醒度の上昇や能動的情報処理と関連し，
一方で，
α帯域およびθ帯域は，
安静状態や抑制的処理，
あるいは低負荷状態と関係すると考えられている．
このため，
EIは認知的負荷や集中度の指標として解釈される．

\subsection{EIにおける集中度指標（CC値）の定義}

本研究では，次式をEngagement Indexによって正規化された
集中度指標CC(ConcentrationCalculation)として定義する．

\begin{equation}
CC = \mathrm{normalize}
\left(
\frac{\beta}{\alpha + \theta}
\right)
\end{equation}

CC値が高い場合，
β帯域活動が相対的に優位であり，
被験者が課題に対して
高い集中状態にあることを示す．

\subsection{EIにおけるリラックス指標（RC値）の定義}

α帯域が優位な状態は，精神的に安定し，
過度な緊張が少ない状態を反映すると考えられている．
本研究では，
次式によりEngagement Indexにおけるリラックス指標RC(RelaxCalculation)値を定義する．

\begin{equation}
RC = \mathrm{normalize}
\left(
\frac{\alpha}{\beta + \theta}
\right)
\end{equation}

RC値が高い状態は，
覚醒度が過度に上昇しておらず，
比較的落ち着いた心理状態を示すものと解釈される．

\subsection{EIにおけるストレス指標（SC値）の定義}

β帯域活動の過度な優位は，
精神的緊張やストレス状態と関連することが知られている．
本研究では，
β帯域とα帯域の比を用いて，
Engagement Indexにおけるストレス指標SC(StressCalculation)値を次式で定義する．

\begin{equation}
SC = \mathrm{normalize}
\left(
\frac{\beta}{\alpha}
\right)
\end{equation}

SC値が高い場合，
β帯域活動が顕著であり，
精神的負荷や緊張状態が高い可能性を示す．

\subsection{正規化処理}

各指標は，
被験者間および時間変化の比較を容易にするため，
事前に設定した最小値および最大値を用いて
0〜100の範囲に正規化している．
この正規化処理により，
CC値，
RC値，
SC値は
直感的に解釈可能な心理状態指標として扱われる．

\subsection{心理量としての解釈}

EIは絶対的な心理状態を直接表すものではなく，
\textbf{同一被験者内での相対的変化量}として解釈される指標である．
特に，安静時（baseline）と課題遂行時のEI差分を用いることで，
課題への関与度の変化を定量的に評価することが可能となる．

本研究では，以下の差分量を用いて解析を行う．

\section{Frontal Alpha/Beta Asymmetry（FAA派生指標）}

\subsection{Frontal Asymmetryの概念}

前頭部脳波の左右非対称性（Frontal Asymmetry）は，
感情処理，動機づけ，注意制御などの心理状態と密接に関係することが
示されている．
特に前頭部$\alpha$波の左右差は，
接近動機と回避動機の指標として広く用いられてきた．

一般に，$\alpha$波パワーは皮質活動の抑制度合いを反映すると考えられており，
$\alpha$波パワーが低いほど局所的な神経活動が高い状態を示す．
そのため，左右前頭部の$\alpha$波パワー差を用いることで，
機能的な活動非対称性を定量化することが可能である．

\subsection*{接近動機と回避動機}

接近動機（approach motivation）とは，
報酬獲得や目標達成といった
正の結果に向かって行動を選択しようとする動機づけを指す．
一方，
回避動機（avoidance motivation）は，
罰や失敗， 不快刺激を回避することを目的とした
行動選択に関与する動機づけである．

接近動機は主に左前頭部の活動増加と関連し，
回避動機は右前頭部の活動増加と関連すると報告されている．
脳波研究では，
前頭部α帯域のパワー低下は局所的な皮質活動の亢進を反映すると考えられており，
左右前頭部α帯域の非対称性は，
接近・回避傾向を推定する指標として広く用いられてきた．

このような前頭部活動の側性は，
単なる情動価の違いにとどまらず，
課題に対する能動的関与や行動選択の方向性とも関連する．
したがって，
本研究で用いるFrontal Alpha/Beta Asymmetry派生指標は，
被験者が課題に対して
接近的に関与しているか，
あるいは回避的に処理しているかを
反映する可能性がある．

\subsection{左右前頭部活動に基づく心理指標の定義}

本研究では，前頭部左右差および帯域別パワー分布に基づき，
集中度（CC），リラックス度（RC），ストレス度（SC）を表す
3つの心理指標を定義する．

\subsection{FAAにおける集中度指標（CC値）の定義}

集中度指標CCは，
左右前頭部における$\beta$帯域活動の非対称性に基づいて定義される．
$\beta$波は注意集中や課題遂行時の能動的認知処理と関連することが知られており，
特に左前頭部の$\beta$活動増加は，
目標志向的・能動的な情報処理状態を反映するとされる．

本研究では，次式によりCCを算出する．

\begin{equation}
CC
=
\log \left( P_{\beta}^{\text{left}} \right)
-
\log \left( P_{\beta}^{\text{right}} \right)
\end{equation}

ここで，
$P_{\beta}^{\text{left}}$ および $P_{\beta}^{\text{right}}$ は，
それぞれ左前頭部（AF7），右前頭部（AF8）における
$\beta$帯域パワーを表す．

CCが正の値をとる場合，
左前頭部$\beta$活動が優位であり，
高い集中状態または能動的思考状態にあると解釈される．

\subsection{FAAにおけるリラックス度指標（RC: Relaxation Component）}

リラックス度指標RCは，
左右前頭部における$\alpha$帯域パワーの平均値を用いて定義される．
$\alpha$波は覚醒水準の低下や抑制的状態と関連し，
精神的な安静状態やリラックス状態の指標として広く用いられている．

RCは次式で与えられる．

\begin{equation}
RC
=
\frac{
P_{\alpha}^{\text{left}}
+
P_{\alpha}^{\text{right}}
}{2}
\end{equation}

RCが高い値を示す場合，
前頭部全体において$\alpha$活動が強く，
低負荷で安静な心理状態にあることを示唆する．

\subsection{FAAにおけるストレス度指標（SC: Stress Component）}

ストレス度指標SCは，
左右前頭部における$\alpha$帯域活動の非対称性に基づいて定義される．
一般に，右前頭部の相対的活動増加は，
回避動機や不安，ストレス反応と関連することが報告されている．

本研究では，次式によりSCを算出する．

\begin{equation}
SC
=
\log \left( P_{\alpha}^{\text{right}} \right)
-
\log \left( P_{\alpha}^{\text{left}} \right)
\end{equation}

SCが正の値をとる場合，
右前頭部$\alpha$活動が優位であり，
ストレス負荷や回避的心理状態の増大を示すと解釈される．

\subsubsection{正規化処理}

各指標は被験者間および時間変化の比較を容易にするため，
事前に設定した範囲に基づき0〜100のスケールへ正規化される．
この処理により，
各指標は直感的かつ相対的な心理量として扱うことが可能となる．


\subsection{本研究における位置づけ}

Frontal Alpha/Beta Asymmetry指標は，
左右非対称性という空間的情報を利用する点において，
周波数比に基づくEngagement Indexとは異なる特徴を持つ．

\section{Sample Entropyに基づく複雑度指標}

\subsection{脳波の複雑性と非線形指標}

脳波信号は非定常かつ非線形な時系列データであり，
単純な振幅や周波数成分のみでは，
認知状態の変化を十分に捉えられない場合がある．
特に注意制御や課題遂行時には，
信号の規則性や不規則性といった「複雑性」が
重要な情報を含むと考えられている．

このような背景から，
時系列の自己相似性に基づく非線形指標として
Sample Entropy（SampEn）が広く用いられている．

\subsection{Sample Entropyに基づくCC・RC・SC算出式の定義}

Sample Entropyは，時系列信号におけるパターンの不規則性を表す指標であり，
値が大きいほど信号が複雑で予測困難であることを意味する．
本研究では，$\beta$帯域（13--30 Hz）にフィルタリングした脳波信号に対して
Sample Entropyを算出し，
認知負荷や覚醒状態の変化を反映させることを目的とする．

\subsection{Sample Entropyの算出}

Sample Entropyは以下の手順で計算される．
埋め込み次元 $m$ の信号パターンが，
次元 $m+1$ に拡張された場合にも類似性を保つ確率を評価し，
その負の対数を取ることで定義される．

\begin{equation}
\mathrm{SampEn} = -\log \left( \frac{A}{B} \right)
\end{equation}

ここで，
$B$ は長さ $m$ の部分系列間で距離が許容範囲 $r$ 以下となる組の数，
$A$ は長さ $m+1$ の部分系列間で同様の条件を満たす組の数である．
Sample Entropyが大きいほど信号の複雑性が高いことを示す．

\subsection{SampEnにおける集中度指標（CC値）の定義}

集中度 CC は，
Sample Entropyそのものを正規化することで定義される．
本研究では，$\beta$帯域の複雑性増加が
注意集中や認知的処理の活性化と対応すると仮定し，
Sample Entropyが大きいほど集中状態が高いと解釈する．

\begin{equation}
CC = \mathrm{normalize}(\mathrm{SampEn}， S_{\min}， S_{\max})
\end{equation}

ここで，$S_{\min}$ および $S_{\max}$ は
実験的に定めたSample Entropyの最小値および最大値であり，
$0$--$100$の範囲に線形正規化を行う．
CCは，課題遂行時の能動的思考や注意資源の投入度を表す指標である．

\subsection{SampEnにおけるリラックス度指標（RC値）の定義 }

リラックス度 RC は，
集中度とは逆の概念として，
Sample Entropyの低下を指標化した量である．
Sample Entropyが小さい状態は，
脳波信号がより規則的で安定していることを意味し，
安静状態やリラックス状態と対応すると考えられる．

\begin{equation}
RC = \mathrm{normalize}(1 - \mathrm{SampEn}， 1 - S_{\max}， 1 - S_{\min})
\end{equation}

この定義により，
RCはCCと補完的な関係を持ち，
高いRC値は低覚醒・低認知負荷状態を示す．

\subsection{SampEnにおけるストレス度指標 （SC値）の定義}

ストレス度 SC は，
Sample Entropy計算過程における
$m$ 次元と $m+1$ 次元のパターン一致率の比に基づいて定義される．
本研究では，
$m+1$ 次元への拡張によって急激に一致率が低下する状態を，
不安定かつストレスの高い状態と解釈する．

\begin{equation}
SC = \mathrm{normalize}\left(
\log \left( \frac{A}{B} \right)，
L_{\min}， L_{\max}
\right)
\end{equation}

ここで，$\log(A/B)$ は
Sample Entropy算出時に用いられる内部比率を表し，
$L_{\min}$，$L_{\max}$ は正規化のための下限・上限値である．
SCが高い値を示す場合，
神経活動が不安定化し，
精神的緊張やストレス状態にある可能性が高いと解釈される．

\section{H30年度卒業生の研究に基づく経験的CC・RC・SC算出式}

本研究では，
理論モデルに基づく指標に加え，
過去の脳波応用研究において経験的に考案された
心理量算出式も比較対象として用いる．
これらの式は，
厳密な神経生理モデルに基づくものではないものの，
$\alpha$帯域および$\beta$帯域パワーの相対関係を利用し，
直感的かつ安定した心理状態推定を可能にすることを目的として設計されている．

以下では，
これらの既存手法に基づく
集中度（CC），
リラックス度（RC），
ストレス度（SC）の算出式について説明する．

\subsection{集中度指標 CC の定義}

集中度 CC は，
$\beta$帯域活動の増加と，
$\alpha$帯域活動の低下が
集中状態を反映するという仮定に基づいて定義される．
本研究で用いる既存式は次式で表される．

\begin{equation}
CC=\min 
\{100， \lfloor\frac{\beta}{2}(1+\frac{1}{\alpha})
\times\frac{100}{2}\rfloor\}
\nonumber \\ 
\end{equation}

この式では，
$\beta$帯域の強度を主要因としつつ，
$\alpha$帯域が低いほど値が増加するよう設計されている．
すなわち，
覚醒度が高く抑制成分が弱い状態を
高い集中状態として評価する構造となっている．
算出結果は0〜100の範囲にクリッピングされ，
集中度指標として扱われる．

\subsection{リラックス度指標 RC の定義}

リラックス度 RC は，
$\beta$帯域活動の低下と
$\alpha$帯域活動の増加を
リラックス状態の特徴と仮定して定義される．
既存研究に基づく算出式は以下の通りである．

\begin{equation}
RC=\min 
\{100， \lfloor(\max \{0，(1．0-\frac{\beta}{3})\}+\frac{\alpha}{2})
\times\frac{100}{2}\rfloor\} 
\nonumber \\
\end{equation}

この式では，
$\beta$帯域が一定値以下である場合に正の寄与を持ち，
過度な覚醒状態を抑制する構造を持つ．
同時に，
$\alpha$帯域活動が高いほど
RC値が増加するよう設計されており，
精神的に安定した安静状態を反映する指標となっている．

\subsection{ストレス度指標 SC の定義}

ストレス度 SC は，
$\alpha$帯域の低下と
$\beta$帯域の相対的増加が
精神的緊張やストレス状態と対応するという
経験的知見に基づいて定義される．
算出式は以下の通りである．

\begin{equation}
SC=\min 
\{100，\lfloor(\max \{0，\frac{1．0-\frac{\alpha}{3}}{5}\}+(\frac{\frac{\beta}{2\alpha}\times4}{5}))
\times100 \rfloor\} 
\nonumber
\end{equation}

この式では，
$\alpha$帯域の減少による抑制解除成分と，
$\beta/\alpha$比による覚醒・緊張成分を
重み付けして合成している．
結果として，
$\beta$帯域が優位で，
かつ$\alpha$帯域が低下した状態において
高いSC値を示す構造となっている．

\section{各心理量算出手法の比較と特徴}

本研究では，
脳波に基づく心理状態推定手法として，
以下の4種類の算出式を用いた．

\begin{itemize}
\item 既存研究に基づく経験的算出式
\item Engagement Index（EI）ベース式
\item Frontal Alpha/Beta Asymmetry（FAA派生）式
\item Sample Entropy（複雑度）ベース式
\end{itemize}

これらは，
理論的背景，
利用する脳波特徴量，
および心理量との対応関係が異なる．
本節では，
各手法の特徴を整理し，
それぞれの利点および限界について比較する．

\subsection{各算出手法の基本的特徴}

既存研究に基づく経験的算出式は，
$\alpha$帯域および$\beta$帯域の相対的関係を
重み付きで組み合わせたものであり，
直感的で安定した数値が得られる点が特徴である．
一方で，
理論的根拠が明示的でなく，
神経生理学的解釈には慎重さが求められる．

Engagement Index（EI）ベース式は，
$\beta/(\alpha+\theta)$ という
明確な理論式を基盤としており，
集中度を表す指標として多くの先行研究で用いられている．
RCおよびSCは，
このEIの思想を拡張する形で定義されているため，
理論的一貫性が高い．

FAA派生式は，
左右前頭部の非対称性に着目した手法であり，
接近・回避動機や情動価との関連が示唆されている．
空間的情報を活用できる点が強みであるが，
電極配置や個人差の影響を受けやすい．

Sample Entropyベース式は，
周波数成分ではなく
脳波信号の時間的複雑性に基づく手法であり，
非線形的な脳活動変化を捉えられる点が特徴である．
一方で，
心理量との対応付けは間接的であり，
解釈には補助的説明が必要となる．

\subsection{手法間の比較表}

各算出手法の特徴を，
利用情報，理論背景，メリット，デメリットの観点から
表\ref{tab:method_comparison}にまとめる．

\begin{table}[h]
\centering
\caption{心理量算出手法の比較}
\label{tab:method_comparison}
\begin{tabular}{|p{3cm}|p{3cm}|p{4cm}|p{4cm}|}
\hline
\textbf{手法} & \textbf{主な特徴量} & \textbf{メリット} & \textbf{デメリット} \\ \hline
既存経験式 &
$\alpha，\beta$帯域パワー &
実装が容易． 数値が安定． 直感的 &
理論的根拠が弱い． 汎用性に課題 \\ \hline
EIベース式 &
$\beta/(\alpha+\theta)$ &
理論的背景が明確． 集中度評価に強い &
情動要素の表現は限定的 \\ \hline
FAA派生式 &
左右前頭部の$\alpha，\beta$非対称性 &
接近・回避動機を反映可能． 空間情報を活用 &
電極配置や個人差の影響が大きい \\ \hline
Sample Entropy式 &時間的複雑度 &非線形的変化を捉えられる． 周波数非依存 &心理量との対応が間接的． 解釈が難しい \\ \hline
\end{tabular}
\end{table}

\subsection{CC・RC・SCによる共通評価軸の意義}

本研究では，
異なる理論背景を持つ手法を
CC（集中），RC（リラックス），SC（ストレス）
という共通の心理量軸に変換することで，
手法間の比較を可能とした．
この枠組みにより，
各手法が
「どの心理側面に敏感であるか」
「課題による変化をどの程度捉えられるか」
を定量的に評価できる．

\subsection{総合的考察}

以上より，
単一の算出式によって
すべての心理状態を正確に表現することは困難であるが，
複数の手法を併用し，
共通指標で比較することで，
より多面的な認知・情動状態の評価が可能となる．
本研究の枠組みは，
用途や目的に応じた指標選択の指針を与えるものである．

\subsection{心理量としての位置づけ}

以上の定義により，
CCは能動的集中状態，
RCは安静・リラックス状態，
SCは緊張・ストレス状態をそれぞれ表す指標として整理される．
これらは絶対値ではなく，
ベースラインとの差分や時間変化を中心に解釈することで，
被験者の認知状態の推移を評価する．


\subsection{心理状態評価への応用}

複雑度指標は，
覚醒度や感情価といった一次元的心理量ではなく，
認知処理の柔軟性や情報統合状態を反映する指標として解釈される．

本研究では，
安静時と課題遂行時の差分量を用い，
以下の指標を算出する．

\begin{equation}
\Delta CC = CC^{\text{task}} - CC^{\text{baseline}}
\end{equation}

\begin{equation}
\Delta RC = RC^{\text{task}} - RC^{\text{baseline}}
\end{equation}

\begin{equation}
\Delta SC = SC^{\text{task}} - SC^{\text{baseline}}
\end{equation}

これらの差分は，
課題遂行に伴う脳波構造の
規則性・持続性・複雑性の変化を表すものである．


 \chapter{脳波計測装置 Muse S}
本章では本研究で使用するMuse Sと呼ばれる脳波計測装置の概要．及び操作方法の説明を行う．
脳波計を利用した研究には人間の思考状態に関する研究や快・不快に関するなどが行われており，
過去の卒業生の脳波に関連する研究においても脳波計が使用されていたため，今回の脳状態の推定に脳波計を使用した．
装着が比較的に容易で装着者の行動を制限しない点や装着者にかかる負荷が小さい点、MuseSに対応したライブラリ等が充実している点を鑑みて，
本研究ではMuseSを使用する．

%---------------------------------------------------------------------------------------------------------------- Muse S の基本構成  --------------------------------------------------------------------------------------------------------------%
\section{Muse S の基本構成}
Muse S は， カナダの InteraXon 社によって開発されたヘッドバンド型の脳波計測装置である． 本装置は， 前頭部および左右側頭部に配置された複数の電極を用いて脳波を計測する構成となっている．
Muse S の外観を図4．1に示す．
本装置は軽量であり， 長時間装着しても被験者の負担が小さい設計となっている． また， ドライ電極を採用しているため， 電極ペーストやジェルを使用する必要がなく， 簡便に脳波計測を行うことが可能である．
\begin{figure}
	\centering
		\includegraphics[width=7cm]{./5_Metaheuristic/Image/muses.eps}
		\caption{MuseSの外観}
	\label{formation}
\end{figure}
%---------------------------------------------------------------------------------------------------------------- Muse S の主な機能  --------------------------------------------------------------------------------------------------------------%
\section{Muse S の主な機能}
Muse S は， 脳波計測を中心とした複数の機能を備えたウェアラブルデバイスである．
第一に， 脳波のリアルタイム計測機能が挙げられる． Muse S は， $\delta$波， $\theta$ 波， $\alpha$ 波， $\beta$ 波といった周波数帯域別の脳波成分を抽出し， 一定周期で外部デバイスへ送信することが可能である． これにより， 被験者の覚醒状態や集中状態， リラックス状態の変化を即時に把握できる．
第二に， 生脳波データの取得機能である． 周波数帯域別データに加え， 電極ごとの生 EEG 波形を取得できるため， 後処理による特徴量抽出や解析にも対応可能である．
第三に， 無線通信機能である． Muse S は Bluetooth Low Energy(BLE) を用いた通信方式を採用しており， PC やスマートフォンとワイヤレスで接続できる．

%----------------------------------------------------------------------------------------------------------------  計測可能な脳波データ  --------------------------------------------------------------------------------------------------------------%
\section{計測可能な脳波データ}
Muse S では， 主に以下の周波数帯域に分類された脳波データを取得することが可能である．$\delta$ 波，$\theta$ 波，$\alpha$ 波，$\beta$ 波
これらの周波数帯域は， 被験者の精神状態や覚醒度と密接な関係を有している． $\alpha$ 波は安静時やリラックス状態で増加する傾向があり， $\beta$ 波は集中や緊張状態において顕著となる． $\theta$ 波は瞑想状態や浅い睡眠時に， $\delta$ 波は深い睡眠時に主に観測される．
本研究では， これらの周波数帯域ごとのパワー値および比率を特徴量として利用する．

\section{Musesの機能}
MuseS本体の部分ごとの名称と機能について説明する．\\

1.電源ボタン：Muse Sの電源を入れたり切ったりするためのボタン．一回押しがbluetoothのペアリングモード，長押しが電源のオン・オフとなる．\\
\\
2.充電ポート：Muse Sの充電を行うためのポート．USB Type-Bコネクタを使用している．フル充電で最大10時間の稼働が可能．\\
\\
3.導電部：Muse Sの電極部分であり，頭皮と接触して脳波を計測する．額部分と耳の裏側部分に配置されている．\\
\\
4.ヘッドバンド：Muse Sを頭部に固定するためのバンド部分．伸縮性のある素材でできており，様々な頭のサイズに対応可能．\\

%---------------------------------------------------------------------------------------------------------------- Muse S の装着方法  --------------------------------------------------------------------------------------------------------------%
\section{Muse S の装着方法}
Muse S はヘッドバンド型の構造を採用しており， 比較的容易に装着することが可能である． 正確な脳波計測を行うためには， 電極が適切な位置で皮膚と接触していることが重要となる．
装着方法と装着例図4.2を以下に示す．\\
\\
1. Muse S の充電ポートの電源ボタンを一回押す．ボタンを押すと反対側のLEDがオレンジ色に点灯する．\\
\\
2.装着者の頭部の大きさに合わせてヘッドバンドの長さを調整する．調整はバンドの後頭部左側にあるスライド式のアジャスターを使用する．\\
\\
3.導電部が額に密着するように，装着を行う．皮膚と電極の間に髪の毛が入らないように注意する．\\
\\
4.磁石式の留め具を使用して，ヘッドバンドを装着する．\\

\begin{figure}
	\centering
		\includegraphics[width=10cm]{./5_Metaheuristic/Image/muses2.eps}
	\caption{Muse S の装着例}
	\label{formation}
\end{figure}
%---------------------------------------------------------------------------------------------------------------- Bluetooth 接続方法と BLED112 ドングル  --------------------------------------------------------------------------------------------------------------%

\section{Bluetooth 接続方法と BLED112 ドングル}
Muse S は Bluetooth Low Energy(BLE) による無線通信を用いて外部デバイスと接続する．Andorid や iphoneなどのスマートフォンと接続することも可能であるが，専用のアプリ以外での接続が出来ず， 取得したデータはアプリ上でのみ閲覧が可能でデータ保存が出来ない仕様となっている．
専用のアプリを用いずに，Windows， MacOS， Linux などの PCとBluetooth接続を行うことで，脳波情報を読み取り，データの抽出することが可能である．
PC 上で安定したデータ取得を行うため， 本研究では BLE 通信専用の USB ドングルである BLED112 を使用した．
本稿では，本研究の開発環境であるWindows11の接続方法について説明する．\\

1. BLED112 ドングルを PC の USB ポートに接続する．\\
\\
2.VScode などのターミナルソフトウェアを起動し， BLED112 ドングルの COM ポートを確認する．\\
\\
3. Muse S の電源を入れ， Bluetooth ペアリングモードにする．\\
\\
4. ターミナルソフトウェア上で "musesls stream" コマンドを実行し， Muse S との接続を行う．\\
\\
5. 接続が成功すると， 周波数帯域別の脳波データがリアルタイムにターミナル上に表示される．\\
\\
対応OSと必要スペックについては以下の通りである．

\begin{table}[htbp]
  \centering
  \caption{Muse S の対応OS（Windows）}
  \label{tab:muse_windows}
  \begin{tabular}{|c|c|}
    \hline
    項目 & 内容 \\
    \hline
    対応OS & Windows 10 以降 \\
    \hline
    対応機種 & Bluetooth Low Energy 対応PC \\
    \hline
    CPU & 公式要件なし（OS要件に準拠） \\
    \hline
    メモリ & 公式要件なし（OS要件に準拠） \\
    \hline
    備考 & 公式Muse App非対応（SDK/研究用途のみ） \\
    \hline
  \end{tabular}
\end{table}

\begin{table}[htbp]
  \centering
  \caption{Muse S の対応OS（macOS）}
  \label{tab:muse_macos}
  \begin{tabular}{|c|c|}
    \hline
    項目 & 内容 \\
    \hline
    対応OS & macOS 12 以降 \\
    \hline
    対応機種 & Intel / Apple Silicon（M1以降）Mac \\
    \hline
    CPU & 公式要件なし（OS要件に準拠） \\
    \hline
    メモリ & 公式要件なし（OS要件に準拠） \\
    \hline
    備考 & 公式Muse App非対応（SDK/研究用途のみ） \\
    \hline
  \end{tabular}
\end{table}

\begin{table}[htbp]
  \centering
  \caption{Muse S の対応OS（Linux）}
  \label{tab:muse_linux}
  \begin{tabular}{|c|c|}
    \hline
    項目 & 内容 \\
    \hline
    対応OS & Linux（Ubuntu等，BLE対応環境） \\
    \hline
    対応機種 & Bluetooth Low Energy 対応PC \\
    \hline
    CPU & 公式要件なし（OS要件に準拠） \\
    \hline
    メモリ & 公式要件なし（OS要件に準拠） \\
    \hline
    備考 & 公式Muse App非対応（SDK/研究用途のみ） \\
    \hline
  \end{tabular}
\end{table}

\begin{table}[htbp]
  \centering
  \caption{Muse S の対応OS（iOS）}
  \label{tab:muse_ios}
  \begin{tabular}{|c|c|}
    \hline
    項目 & 内容 \\
    \hline
    対応OS & iOS 15 以上 \\
    \hline
    対応機種 & iPhone / iPad（Bluetooth LE対応） \\
    \hline
    CPU & Apple A12 Bionic 以降（推奨） \\
    \hline
    メモリ & 3 GB 以上（推奨） \\
    \hline
    備考 & 公式Muse App対応（主要利用環境） \\
    \hline
  \end{tabular}
\end{table}

%---------------------------------------------------------------------------------------------------------------- データ取得および保存方法  --------------------------------------------------------------------------------------------------------------%
\section{データ取得および保存方法}
取得した脳波データは， 時系列データとして PC 上に保存される． 本研究では， 一定時間ごとにデータを区切り， 周波数帯域別の平均値や比率を算出することで， メンタル状態推定に用いる特徴量を生成した．
この構成により， 被験者の状態変化をリアルタイムに追跡しつつ， 後処理による解析も可能になっている．

%---------------------------------------------------------------------------------------------------------------- 本研究における Muse S の位置づけ  --------------------------------------------------------------------------------------------------------------%
\section{本研究における Muse S の位置づけ}
本研究において Muse S は， 被験者の内的状態を客観的に取得するための主要なセンサとして位置づけられる． 自己申告やテキスト入力のみに依存せず， 脳波のバイタルデータを用いることで， 潜在的なメンタル状態を推定できる点が本研究の特徴である．
また， 装着負担が小さい Muse S は， 長時間計測や日常環境での利用にも適しており， 将来的な実用システムへの応用可能性も高い．


\chapter{メンタル推定システム}
本章では本研究で考案したメンタル推定システムの全体構成及び各モジュールの詳細について説明する．
脳波を利用したシステム開発の研究として脳波と脳状態との関連性から，その特性を利用した学習のサポートシステムの開発や，ユーザのメンタル状態を推定して能動的なフィードバックを行うシステムの開発などが行われている．
本システムでは，主に「Muse Sにより取得した脳波データのリアルタイム表示」，「取得した脳波データに基づくメンタル状態の推定」を行う．
このシステムを通したEIやFAA，SampEn，卒業生の研究に基づくメンタル推定計算式の各4つの指標の精度評価実験については，第6章で述べる．

\section{メンタル推定システムの概要}
本システムは， VScode上で動作するPythonプログラムとして実装されている．MuseSとPCにBluetooth接続を行い， Muse S から送信される脳波データをリアルタイムに取得する．
取得した脳波データは， 一定時間ごとに区切られ， 周波数帯域別のパワー値や比率を算出することで， メンタル状態推定に用いる特徴量を生成する．
これらの特徴量を基に， 集中度（CC）， リラックス度（RC）， ストレス度（SC）を算出し， 画面上にリアルタイムで表示する．

\section{メンタル推定システムの操作方法}
本システムの操作方法について説明する．MuseSの接続からシステムの実行，及びデータの保存方法について述べる．
\subsection{Muse S の接続}
VScodeを開き，ターミナル上で"musesls stream"コマンドを実行し，Muse SとのBluetooth接続を行う．
接続が成功すると，ターミナルに"steaming EEG"と表示される．

\begin{figure}[htbp]
	\centering
		\includegraphics[width=15cm]{./5_Metaheuristic/Image/setuzoku1.jpg}
	\caption{MuseSの接続画面}
	\label{formation}
\end{figure}

\begin{figure}[htbp]
	\centering
		\includegraphics[width=15cm]{./5_Metaheuristic/Image/setuzoku2.jpg}
	\caption{MuseSの接続画面}
	\label{formation}
\end{figure}

\subsection{システムの実行}
MuseSがPCと接続した後，VScode上でメンタル推定システムのPythonプログラムを実行する。
プログラム実行後，10秒間の安静時データを取得し，ベースラインデータとして保存する．
その後，リアルタイムで脳波データの取得とメンタル状態の推定が開始される．

\begin{figure}[htbp]
	\centering
		\includegraphics[width=15cm]{./5_Metaheuristic/Image/stream1.jpg}
	\caption{平常時ベースライン(10秒)取得画面}
	\label{formation}
\end{figure}

\newpage
\subsection{リアルタイムのメンタル状態の推定}
取得した脳波データに基づき，リアルタイムでメンタル状態の推定を行う．
各周波数帯域のパワー値や比率を算出し，集中度（CC），リラックス度（RC），ストレス度（SC）を計算する．
計算結果は，画面上にリアルタイムで表示される．
計測したデータは，CSV形式で保存され，後処理による解析に利用できる．
CSVファイルの内訳は以下の通りである．
pythonプログラムのコード内の計算式をEI，FAA，SampEn，卒業生の研究に基づくメンタル推定計算式の4つの指標に変更することで，各指標のデータを取得可能である．

\begin{table}[htbp]
\centering
\caption{CC, RC, SC指標および関連変数の定義}
\label{tab:ccrcsc_definition}
\begin{tabular}{lll}
\hline
\textbf{列名} & \textbf{意味} & \textbf{用途} \\
\hline
timestamp &
データを記録した時刻 &
実時間との同期 \\

elapsed\_sec &
計測開始からの経過時間 [秒] &
時系列解析 \\

CC &
各時刻で算出された集中指標（Concentration Calculation） &
瞬時的な集中度評価 \\

RC &
各時刻で算出されたリラックス指標（Relax Calculation） &
瞬時的なリラックス度評価 \\

SC &
各時刻で算出されたストレス指標（Stress Calculation） &
瞬時的なストレス度評価 \\

CC\_avg &
CCの移動平均値 &
集中状態の推移把握 \\

RC\_avg &
RCの移動平均値 &
リラックス状態の推移把握 \\
SC\_avg &
SCの移動平均値 &
ストレス状態の推移把握 \\

$\Delta$CC &
CCの変化量（前時刻または平均値との差分） &
集中状態変化の検出 \\

$\Delta$RC &
RCの変化量（前時刻または平均値との差分） &
リラックス状態変化の検出 \\

$\Delta$SC &
SCの変化量（前時刻または平均値との差分） &
ストレス状態変化の検出 \\
\hline
\end{tabular}
\end{table}

\newpage

\begin{figure}[htbp]
	\centering
		\includegraphics[width=11cm]{./5_Metaheuristic/Image/stream3.jpg}
	\caption{各周波数のリアルタイム表示画面}
	\label{formation}
\end{figure}

\begin{figure}[htbp]
	\centering
		\includegraphics[width=11cm]{./5_Metaheuristic/Image/stream2.jpg}
	\caption{メンタル推定システムのリアルタイム表示画面}
	\label{formation}
\end{figure}


\chapter{考案システムの評価実験}

本章では，
本研究で提案した脳状態推定システムの有効性を検証するために実施した
評価実験について述べる．
本システムでは，
脳波信号から集中度（CC），リラックス度（RC），ストレス度（SC）を推定する複数の指標を導入しており，
それらの精度および特性を比較することが重要となる．

本章前半では，
既存研究および本研究で採用した複数の心理量算出指標について，
同一実験条件下での精度比較実験を行い，
各指標の特徴と有効性を検証する．
後半では，
過去の研究で提案された経験的算出式に対し，
本研究独自にパラメータ調整を施した新しい算出式を提案し，
その精度向上の有無を評価する検証実験について述べる．

\section{各指標の精度比較実験}

\subsection{実験目的}

本実験の目的は，
脳波から心理状態を推定する複数の指標について，
課題負荷による変化をどの程度適切に捉えられるかを比較し，
本システムにおいて有効な指標を明らかにすることである．

比較対象とした指標は以下の4種類である．

\begin{itemize}
\item 既存研究に基づく経験的CC・RC・SC算出式
\item Engagement Index（EI）ベース指標
\item Frontal Alpha/Beta Asymmetry（FAA派生）指標
\item Sample Entropy（複雑度）ベース指標
\end{itemize}

これらは理論的背景および利用する特徴量が異なるため，
同一の実験条件において比較することで，
それぞれの特性を明確にできると考えられる．

\subsection{被験者}

被験者は健常な大学生・大学院生○名（男性○名，女性○名）とした．
すべての被験者に対して事前に実験内容を説明し，
インフォームドコンセントを得た上で実験を実施した．
被験者はいずれも神経疾患や精神疾患の既往を有していない．

\subsection{実験環境}

脳波計測には，
前頭部計測が可能な簡易型脳波計を用いた．
電極配置はAF7およびAF8とし，
左右前頭部の活動を同時に取得した．
サンプリング周波数は256 Hzとした．

実験は静かな室内環境で行い，
被験者は椅子に座った状態で計測を行った．

\begin{figure}[h]
\centering
\includegraphics[width=0.7\linewidth]{experiment_setup.png}
\caption{実験環境および脳波計測の様子}
\label{fig:experiment_setup}
\end{figure}

\subsection{実験手順}

実験は，
ベースライン計測区間と課題遂行区間から構成される．

まず，
被験者に対して安静状態を保つよう指示し，
最初の30秒間をベースライン区間として脳波を計測した．
この区間では，
画面注視のみを行い，
認知的負荷を与えない状態とした．

次に，
被験者に対して課題を提示し，
約60秒間の課題遂行区間を設けた．
課題内容は，
注意集中および能動的思考を要する作業とし，
ベースライン区間との差が明確に生じるよう設計した．

\begin{figure}[h]
\centering
\includegraphics[width=0.9\linewidth]{experiment_timeline.png}
\caption{実験手順のタイムライン}
\label{fig:experiment_timeline}
\end{figure}

\subsection{評価指標}

各算出式により得られたCC，RC，SCについて，
以下の観点から評価を行った．

\begin{itemize}
\item ベースライン区間と課題区間の平均値の差（Δ値）
\item 被験者間でのばらつき
\item 指標ごとの感度の違い
\end{itemize}

特に，
課題負荷による変化量が大きい指標ほど，
心理状態推定において有効であると仮定した．

\subsection{実験結果}

各指標について，
ベースライン区間および課題区間の平均値を算出し，
その差分を求めた．
結果の一例を図\ref{fig:delta_comparison}に示す．

\begin{figure}[h]
\centering
\includegraphics[width=0.9\linewidth]{delta_comparison.png}
\caption{各指標におけるΔ値の比較}
\label{fig:delta_comparison}
\end{figure}

結果より，
EIベース指標およびFAA派生指標では，
CC値において課題区間での明確な増加が確認された．
一方，
Sample Entropyベース指標では，
SC値の変化が比較的小さく，
課題負荷への感度にばらつきが見られた．

\subsection{考察}

これらの結果から，
EIベース指標は集中状態の変化を安定して捉える能力が高いことが示唆される．
FAA派生指標は，
左右前頭部の活動差を反映するため，
個人差の影響を受けやすい一方で，
接近的認知状態の検出に有効であると考えられる．

Sample Entropyベース指標は，
非線形的特徴を捉えられる利点を有するが，
心理量との対応付けには慎重な解釈が必要である．

\section{パラメータ調整式の精度測定実験}

\subsection{実験目的}

本節では，
過去研究で提案された経験的算出式に対し，
本研究でパラメータ調整を施した新しい算出式の有効性を検証する．
特に，
課題負荷に対する感度向上が得られるかを評価する．

\subsection{パラメータ調整方法}

既存式における係数および正規化範囲について，
予備実験データを用いて調整を行った．
調整の方針として，
ベースライン区間と課題区間の差が最大化されるよう，
各パラメータを探索した．

\subsection{評価方法}

調整前後の算出式について，
同一被験者・同一実験データを用いてCC，RC，SCを算出し，
Δ値の比較を行った．

\subsection{結果および考察}

調整後の算出式では，
特にCC値において課題区間での増加量が大きくなり，
集中状態の検出感度が向上した．
一方で，
RCおよびSCについては，
過度な変動を抑制する必要性も確認された．

以上より，
パラメータ調整は指標の感度向上に有効であるが，
安定性とのトレードオフを考慮する必要があることが示された．

%%%%%%%%%%%%%%%%%%%%%%%%%%%%%%%%%%%%%%%%%%


%%%%%%%%%%%%%%%%%%%%%%%%%%%%%%%%%%%%%%%%%%%
\expandafter\ifx\csname ifdraft\endcsname\relax
  \end{document}
\fi
